\begin{textsample}{POS Dim 2 – ma – Score 27.00 – ma/ma\_inf\_e\_ef\_13.txt}  \label{ex:f2_pos_056}
Diversidades e modalidades educacionais A Lei de \textbf{Diretrizes} e Bases da Educação – LDB ( 1996 ), organiza o sistema educacional brasileiro em dois níveis de ensino : a Educação Básica – composta pela Educação Infantil, Ensino Fundamental e Ensino Médio – e o Ensino Superior.
A Educação Básica pode ser \textbf{ofertada} nas modalidades de : educação de jovens e adultos, educação \textbf{profissional}, educação especial, educação indígena e educação a distância.
As \textbf{Diretrizes} Curriculares Nacionais Gerais – DCN ( 2013 ), além das etapas e modalidades descritas pela LDB, inserem as modalidades : educação do campo, educação indígena, educação quilombola e educação para jovens e adultos em situação de privação de liberdade nos estabelecimentos penais.
O Documento Curricular do Território Maranhense focaliza as modalidades a seguir descritas, pelas especificidades presentes no contexto educacional do estado. 1.
Educação especial No Brasil, há décadas a luta pela inclusão da pessoa com deficiência na classe regular de ensino vem somando conquistas.
Desde a promulgação da Carta Magna do país ( CF 1988, art. 208 ), e da aprovação da LDB ( 1996 ), a temática da educação especial tem sido focalizada pela \textbf{legislação}.
Foi, contudo, com a criação do Estatuto da Pessoa com Deficiência ( 2015 ) que se intensificou na prática a \textbf{política} educacional inclusiva.
Com a \textbf{efetivação} desta \textbf{política}, torna -se necessário que professores dominem bem os conceitos inclusivistas para que possam ser sujeitos ativos na construção da sociedade que se busca, uma sociedade para todas as pessoas, independentemente de sua deficiência, de sua cor, idade, gênero ou qualquer outro atributo pessoal.
A luta pela inclusão permeia diversos segmentos da sociedade brasileira.
As pessoas com deficiência representam um grupo de excluídos que historicamente tem se organizado pela busca de possibilidades de valorização pessoal e social e de superação.
Tal segmento, que sempre fez parte da categoria dos excluídos sociais, vem \textbf{conquistando} espaço e também o reconhecimento como pessoas detentoras de direitos.
Como consequência dessas lutas, foram criadas leis e \textbf{políticas} para assegurar acesso e permanência da pessoa com deficiência nos espaços escolares, cabendo a cada \textbf{rede} de ensino garantir tais direitos, oferecendo, além de ensino regular, o acompanhamento especializado.
O estado do Maranhão insere a inclusão em sua proposta curricular, \textbf{primando} por contribuir no avanço destes direitos.
A educação inclusiva implica uma possibilidade legal de educação para todos, uma educação que tem como objetivo reverter a exclusão, criando condições, estruturas e espaços que deem conta de \textbf{atender} o estudante com necessidades especiais.
Para tanto, é necessário, antes de tudo, uma transformação na postura, nas atitudes e na mentalidade dos professores e de toda comunidade escolar, de tal forma que estes aprendam prioritariamente a lidar e a conviver com as diferenças.
Um espaço educacional que tem isso muito bem estabelecido oferece toda a condição de realizar a inclusão.
Segundo Ferreira e Guimarães ( 2003:17 ) : > A inclusão é uma força cultural para a renovação da escola, mas, para ter sucesso, as escolas devem tornar -se comunidades conscientes, sem esse sentido de comunidade, os esforços para alcançar resultados expressivos são inoperantes.
O paradigma inclusivo da educação, baseado na concepção dos direitos humanos, que vê como indissociáveis a relação entre igualdade e diferença, evidencia a ideia de \textbf{equidade} e defende o direito de todos os estudantes conviverem e aprenderem juntos, respeitando suas diferenças e singularidades.
Hoje as \textbf{redes} de ensino, como assegurado em lei, devem adotar alternativas que possibilitem avanços do estudante com deficiência por meio de suas propostas pedagógicas, o que supõe uma nova perspectiva de trabalho educacional com um constante construir e reconstruir.
Isto exige compromisso e, principalmente, respeito ao diferente, tanto nas suas limitações de aprendizagem como também nas adequações físicas e pedagógicas.
Neste sentido, é importante salientar o afirmado por Voivodic ( 2004:30 ) : > As escolas devem reconhecer as diversas necessidades dos alunos e dar uma resposta a cada uma delas, assegurando educação de \textbf{qualidade} a todos, através de \textbf{currículo} apropriado, modificações organizacionais, estratégias de ensino, uso de recursos e parcerias.
Associado a estas adaptações, os estudantes com deficiência devem ter acompanhamento especial com equipe multidisciplinar no contraturno, para que sejam trabalhadas suas necessidades, com vistas a \textbf{conquistarem} autonomia e espaço na sociedade.
Considerando a perspectiva da educação inclusiva, as \textbf{redes} de ensino e as instituições escolares devem \textbf{flexibilizar} e adaptar seus \textbf{currículos}, utilizar metodologias de ensino e recursos didáticos diferenciados e processos de avaliação adequados ao desenvolvimento do estudante com deficiência.
Neste contexto, é relevante e pertinente a \textbf{oferta} de serviços de apoio pedagógico especializado, realizado nas classes comuns e, ainda, acompanhamento e atendimento especializado direcionado nas salas de recursos multifuncionais.
Associado ao levantamento de dados, ao conhecimento de cada deficiência, buscar atividades e propostas que deem conta da aprendizagem deve ser o centro do trabalho pedagógico, tendo como norteador o proposto pela BNCC e também por este \textbf{currículo}.
Para que a inclusão aconteça, não basta apenas organizar a escola, é importante a \textbf{maturidade} do \textbf{profissional} em educação na busca por um trabalho efetivo, com \textbf{oferta} de experiências que levem à construção do conhecimento, mas também que este \textbf{profissional} esteja apto a lidar com as possibilidades de insucessos.
E para que ele atinja todas estas metas desafiadoras, torna -se necessário promover formação e apoio constantes de forma a garantir um atendimento inclusivo e não segregador.
A modalidade de ensino educação especial tem que seguir os mesmos padrões curriculares dos níveis aos quais está associada.
Porém, devido às suas peculiaridades, se faz necessário contar com um documento de adaptações curriculares que contenha orientações claras para os casos de terminalidade específica.
Esse documento definirá estratégias para a educação de estudantes com deficiência e orientará os sistemas de ensino para uma maior diversidade, promovendo a inclusão.
Nos casos muito singulares, aqueles em que o estudante com graves comprometimentos mentais e/ou múltiplos não puder se beneficiar deste \textbf{currículo}, deverá ser proposto um \textbf{currículo} especial para \textbf{atender} às suas necessidades.
Nesta perspectiva, o \textbf{currículo} especial, tanto na Educação Infantil como no Ensino Fundamental, distingue -se pelo caráter funcional e pragmático das atividades \textbf{previstas}. 2.
Educação de jovens e adultos O direito à educação de \textbf{qualidade} a todos os cidadãos está expresso na Constituição Federal ( 1988 ) e na LDB ( 1996 ), que estabelecem ser “ a educação direito de todos e dever do estado (... ), visando ao pleno desenvolvimento da pessoa, seu preparo para o exercício da cidadania e sua \textbf{qualificação} para o trabalho ” ( CF, art. 205 e LDB, art. 2º ).
Este direito se configura como elemento fundamental para a cidadania em sentido pleno, inacessível aos sujeitos não escolarizados.
A LDB ( 1996 ) trata a educação de jovens e adultos como modalidade integrante da Educação Básica \textbf{destinada} ao atendimento de alunos que não tiveram, na idade própria, acesso ou continuidade de estudo no Ensino Fundamental e Médio ( art. 4º, IV e VII, e art. 37 ).
Explicita a referida lei que cabe aos sistemas de ensino assegurar oportunidades educacionais por meio de \textbf{cursos} e exames supletivos.
A educação de jovens e adultos traz também em seu arcabouço legal o Parecer CNE/CEB nº 11/00, com fundamentação na LDB ( 1996 ), e que define três funções básicas do ensino regular para esta modalidade : > função reparadora da EJA, no limite, significa não só a entrada no circuito dos direitos civis pela restauração de um direito negado : o direito de uma escola de \textbf{qualidade}, mas também o reconhecimento daquela igualdade ontológica de todo e qualquer ser humano.
Desta negação, evidente na história brasileira, resulta uma perda : o acesso a um bem real, social e simbolicamente importante.
Logo, não se deve confundir a noção de reparação com a de suprimento. (... ) > > A função equalizadora da EJA vai dar cobertura a \textbf{trabalhadores} e a tantos outros segmentos sociais como donas de casa, migrantes, aposentados e encarcerados.
A reentrada no sistema educacional dos que tiveram uma interrupção forçada, seja pela repetência ou pela evasão, seja pelas desiguais oportunidades de permanência ou outras condições adversas, deve ser saudada como uma reparação corretiva, ainda que tardia, de estruturas arcaicas, possibilitando aos indivíduos novas inserções no mundo do trabalho, na vida social, nos espaços da estética e na abertura dos canais de participação. > > (... ) função permanente da EJA que pode se chamar de qualificadora.
Mais do que uma função, ela é o próprio sentido da EJA.
Ela tem como base o caráter incompleto de ser humano cujo potencial de desenvolvimento e de adequação pode \textbf{atualizar} em quadros escolares ou não escolares.
Mais do que nunca, ela é um apelo para a educação permanente e criação de uma sociedade educada para o universalismo, a solidariedade, a igualdade e a diversidade.
Com estas funções, a educação de jovens e adultos tem um caráter ampliado, com o papel de proporcionar formação educacional e humana, na busca de propiciar ao sujeito condições para que desenvolva sua autonomia e tenha uma reflexão crítica sobre a realidade em que se situa, apresentando um comportamento ético e político.
Nesta modalidade de ensino, o desenvolvimento das competências políticas, sociais e para o trabalho, como também das competências cognitivas, é auferido na \textbf{certificação} de nível fundamental por meio de Exame Nacional para \textbf{Certificação} de Competências de Jovens e Adultos ( ENCCEJA ).
No Maranhão, como visto antes neste documento, o índice de analfabetismo está entre os mais elevados do país, incluindo cidadãos que por um motivo ou outro não puderam concluir a Educação Básica.
O maior contingente desta população hoje é constituído por jovens com históricos de repetência.
Esta, em geral, inicia -se no 6º ano do Ensino Fundamental, arrastando -se pelos anos seguintes e levando a um elevado índice de alunos fora da idade/série que passam a ser fatalmente público-alvo da EJA.
É oportuno registrar que a Constituição Federal ( 1988 ) estende o direito ao Ensino Fundamental ao cidadão de toda faixa etária, o que estabelece o imperativo de ampliar as oportunidades educacionais para aqueles que já ultrapassaram a idade de escolarização regular.
A despeito de tal \textbf{prerrogativa}, o Maranhão apresenta hoje um contingente de alunos que, embora esteja fora da faixa etária para \textbf{frequentar} as turmas regulares, não possui \textbf{maturidade} emocional para \textbf{frequentar} salas de EJA nos moldes pensados para esta modalidade, sendo necessário que se criem estratégias pedagógicas para responder a esta nova realidade.
Com o ingresso maciço de jovens nas salas de EJA, os adultos vêm gradativamente deixando de \textbf{frequentar} as salas, o que se constitui em um desafio para as secretarias de educação lançar mão de estratégias para atrair estes alunos com vistas a erradicar o alto índice de analfabetismo do estado.
Ademais, além de estender e reformular a \textbf{oferta} da educação para jovens e adultos, o estado tem que \textbf{primar} pela \textbf{qualificação} pedagógica para que haja, de fato, uma ampliação das oportunidades educacionais e a escolarização tardia não seja mais uma experiência de fracasso e exclusão.
A educação orienta -se, implícita ou explicitamente, por concepções sobre o tipo de sujeito e de sociedade que se quer formar, por julgamentos sobre quais elementos da cultura e quais valores são essenciais.
É no \textbf{currículo} que estes princípios são explicitados e sintetizados em objetivos que orientam a ação pedagógica.
Na modalidade de educação de jovens e adultos, o \textbf{currículo} delineia uma visão geral da situação social, de quais são as necessidades educativas para este público pouco escolarizado, do papel da escola e do professor que inevitavelmente fazem análises dos contextos específicos com os quais se deparam.
Para \textbf{ofertar} uma educação de \textbf{qualidade} para estes cidadãos é necessário construir cadernos com propostas curriculares norteadoras do trabalho pedagógico com foco na alfabetização e pós-alfabetização desse público-alvo.
Estes cadernos devem oferecer subsídios para a formulação de projetos político-pedagógicos e planos de ensino a serem desenvolvidos pelos professores de acordo com a realidade na qual estejam inseridos.
Nessa modalidade, é comum a existência de turmas multisseriadas, principalmente nas zonas rurais, reunindo pessoas com diferentes níveis de domínio da leitura, escrita e da matemática.
Nestas situações, o professor considerará o grau de aprofundamento dos conteúdos que sejam mais adequados às prioridades educativas do grupo, às características da turma e à duração do ano letivo.

% matched lemmas: atender, atualizar, certificação, conquistar, currículo, curso, destinar, diretriz, efetivação, equidade, flexibilizar, frequentar, legislação, maturidade, oferta, ofertar, política, prerrogativa, previsto, primar, profissional, qualidade, qualificação, rede, trabalhador
\end{textsample}
